\documentclass[12pt,a4paper]{article}

% UC3M-style wrapper: cover + courtesy pages + abstract/keywords + body in Markdown.
% Source-of-truth for the body remains in Markdown.

\usepackage{iftex}
\ifPDFTeX
	\usepackage[utf8]{inputenc}
	\usepackage[T1]{fontenc}
	% Times-like fonts para pdfLaTeX.
	% newtxtext/newtxmath funcionan en Overleaf (TeX Live completo).
	% Si compilas localmente con MiKTeX basico, cambia por: \usepackage{mathptmx}
	\usepackage{newtxtext}
	\usepackage{newtxmath}
\else
	% XeLaTeX/LuaLaTeX: allows Times New Roman exactly
	\usepackage{fontspec}
	\setmainfont{Times New Roman}
\fi

% CORRECCION 1: añadir opcion es-tabla para que las tablas digan "Tabla" en vez de "Cuadro"
\usepackage[spanish,es-tabla]{babel}
\usepackage[a4paper,top=2.5cm,bottom=2.5cm,left=3cm,right=3cm]{geometry}
\usepackage[hidelinks]{hyperref}
\usepackage{graphicx}
\usepackage{ragged2e}
\usepackage{setspace}

% CORRECCION 2: paquetes de matematicas necesarios para las formulas $$...$$ del .md
\usepackage{amsmath}
\usepackage{amssymb}

% Render Markdown inside LaTeX
% Enable pipe tables so Markdown tables render correctly.
\usepackage[hybrid,pipeTables,tableCaptions]{markdown}

% Do not render Markdown strong emphasis (**...**) in bold.
% This keeps the body text clean while preserving LaTeX's own heading styles.
\renewcommand{\markdownRendererStrongEmphasis}[1]{#1}

% CORRECCION 3: con inputenc utf8 activo, los caracteres especiales van directos (sin \').
% Usar \' dentro de \renewcommand puede provocar conflictos con babel/spanish.
\renewcommand{\contentsname}{ÍNDICE DE CONTENIDOS}
\renewcommand{\listfigurename}{ÍNDICE DE FIGURAS}
\renewcommand{\listtablename}{ÍNDICE DE TABLAS}

% ---------- Cover fields (fill these) ----------
\newcommand{\tfgDegree}{\textbf{Ingeniería Informática}}
\newcommand{\tfgAcademicYear}{Curso académico: \textbf{2025--2026}}
\newcommand{\tfgTitle}{\textbf{Generación automática de consultas SPARQL desde lenguaje natural}}
\newcommand{\tfgAuthor}{Autor/a: \textbf{Amina Errami Maslaoui}}
\newcommand{\tfgSupervisor}{Tutor/a: \textbf{Jose María Álvarez}}
\newcommand{\tfgPlaceDate}{Lugar y fecha: \textbf{Madrid, abril de 2026}}
% If the work will be published in the UC3M open archive, add the chosen license:
% \newcommand{\tfgLicense}{Licencia: \textbf{Creative Commons (TODO)}}

\begin{document}

% Paragraph layout (justified + 6pt between paragraphs) and line spacing.
\justifying
\setstretch{1.5}
\setlength{\parskip}{6pt}
\setlength{\parindent}{0pt}

% Markdown headings already include manual numbering (e.g., "2.2 ...").
% Disable LaTeX automatic section numbering to avoid duplicated numbers like "0.2.2 2.2".
\setcounter{secnumdepth}{0}

% Preliminaries in Roman numerals (I, II, III, ...)
\pagenumbering{Roman}
\setcounter{page}{1}

% ---------- Cover (Page I) ----------
\begin{titlepage}
\thispagestyle{empty}
\begin{center}
\vspace*{1.5cm}

{\Large \tfgDegree \\
\vspace{0.2cm}
\tfgAcademicYear}

\vspace{2.0cm}
{\LARGE Trabajo de Fin de Grado}\\
\vspace{0.8cm}
{\Huge \tfgTitle}\\

\vfill
{\large \tfgAuthor}\\
\vspace{0.2cm}
{\large \tfgSupervisor}\\
\vspace{0.8cm}
{\large \tfgPlaceDate}\\

% (Optional) License line for open archive publication
% \vspace{0.5cm}
% {\large \tfgLicense}\\

\end{center}
\end{titlepage}

% ---------- Courtesy page (Page II, blank) ----------
\clearpage
\thispagestyle{empty}
\null
\clearpage

% ---------- Resumen + palabras clave (Page III) ----------
\thispagestyle{plain}
\section*{RESUMEN}
% NOTE: El resumen no debe citar bibliografía ni mencionar figuras/tablas.
\noindent
El presente Trabajo de Fin de Grado aborda el diseño e implementación de una herramienta que permite formular preguntas en lenguaje natural sobre un grafo RDF y obtener como salida una consulta SPARQL ejecutable, su resultado y una traza de explicación. La motivación principal es reducir la barrera de entrada de SPARQL en escenarios de auditoría e integridad de datos, donde es necesario consultar información de trazabilidad (requisitos, modelos y evidencias) de forma repetible, verificable y sin depender de servicios externos.

Para alcanzar este objetivo, se propone un pipeline de traducción controlado y determinista que combina (i) normalización del texto de entrada, (ii) indexado del vocabulario observado en el propio grafo, (iii) \emph{grounding} explícito de palabras y expresiones a clases, predicados y operadores, y (iv) compilación de familias de consulta a patrones SPARQL bien delimitados (por ejemplo, ausencia mediante \texttt{FILTER NOT EXISTS} y auditorías de duplicidad mediante \texttt{GROUP BY}/\texttt{HAVING}). Antes de la ejecución, un \emph{checker} valida que la consulta es de solo lectura y que no introduce términos fuera del esquema observado.

La evaluación se realiza mediante la ejecución local de las consultas sobre un grafo de pruebas reproducible y la comparación frente a un conjunto de consultas SPARQL de referencia, además de pruebas de regresión basadas en reformulaciones (paráfrasis). Los resultados ponen de manifiesto que, para el conjunto de intenciones cubiertas por los operadores implementados, el sistema produce consultas ejecutables de forma consistente y aporta una explicación \mbox{(operador seleccionado, señales de \emph{grounding} y patrón de consulta)} que facilita la auditoría y la depuración.

\vspace{0.6cm}
\noindent\textbf{Palabras clave:} SPARQL; RDF; grafos de conocimiento; semantic parsing; KGQA; trazabilidad; auditoría de datos; reproducibilidad.

% ---------- Courtesy page after Resumen (Page IV, blank) ----------
\clearpage
\thispagestyle{empty}
\null
\clearpage

% ---------- Abstract + keywords (Page V) ----------
\thispagestyle{plain}
\section*{ABSTRACT}
\noindent
This Bachelor Thesis presents the design and implementation of a tool that translates natural-language questions into executable SPARQL queries over an RDF graph, returning both query results and an explanation trace. The main objective is to enable non-expert users to perform auditing and traceability checks on knowledge-graph data in a repeatable and verifiable way, while keeping the full pipeline offline and under strict schema control.

The proposed approach is based on a deterministic, operator-driven pipeline that (i) normalizes the input text, (ii) builds an explicit index of the observed graph vocabulary, (iii) grounds relevant tokens and phrases to schema concepts and supported operators, and (iv) compiles each operator into a constrained SPARQL pattern (e.g., absence via \texttt{FILTER NOT EXISTS}, duplicates via \texttt{GROUP BY}/\texttt{HAVING}). Before execution, a dedicated \emph{checker} enforces read-only queries and rejects terms that do not appear in the observed schema.

Evaluation is performed by running the generated queries locally on a reproducible test graph and comparing behavior against reference SPARQL queries, complemented with paraphrase-based regression tests. Results show that, within the set of supported auditing families, the system reliably produces executable queries and provides an explanation trace \mbox{(selected operator, grounding signals, and query pattern)} that supports debugging and auditability.

\vspace{0.6cm}
\noindent\textbf{Keywords:} SPARQL; RDF; knowledge graphs; semantic parsing; KGQA; traceability; data auditing; reproducibility.

% ---------- Courtesy page after Abstract (Page VI, blank) ----------
\clearpage
\thispagestyle{empty}
\null
\clearpage

% ---------- Dedication / Acknowledgements (Page VII) ----------
\thispagestyle{plain}
\section*{}
\vspace*{4cm}
\begin{center}
\textit{(Opcional) Dedicatoria / Agradecimientos.}
\end{center}

% ---------- Courtesy page (Page VIII, blank) ----------
\clearpage
\thispagestyle{empty}
\null
\clearpage

% ---------- TOC / lists (start around Page IX) ----------
\thispagestyle{plain}
\tableofcontents
\clearpage

\thispagestyle{plain}
\listoffigures
\clearpage

\thispagestyle{plain}
\listoftables
\clearpage

% ---------- Abbreviations index ----------
\thispagestyle{plain}
\section*{ÍNDICE DE ABREVIATURAS}
\noindent
\begin{tabular}{@{}ll@{}}
ASK & Consulta booleana en SPARQL \\
CC & Creative Commons \\
CLI & Interfaz de línea de comandos \\
DCTERMS & Dublin Core Metadata Terms \\
FOAF & Friend of a Friend (vocabulario RDF) \\
IRI & Internationalized Resource Identifier \\
KGQA & Question Answering over Knowledge Graphs \\
NB & Naive Bayes \\
NL & Lenguaje natural \\
RDF & Resource Description Framework \\
SPARQL & SPARQL Protocol and RDF Query Language \\
TTL & Turtle (serialización RDF) \\
V\&V & Verificación y Validación \\
W3C & World Wide Web Consortium \\
\end{tabular}

\clearpage

% ---------- Body: Arabic numbering starts at Introduction ----------
\pagenumbering{arabic}
\setcounter{page}{1}

% The file is generated/copied by prepare_overleaf.ps1
\markdownInput{memoria_TFG_esqueleto.md}

\end{document}